\section{Model}
\label{sec:model}

We fix the estimand, the source model, and the holdover state so that the two
regimes---cross-checked and coasting---are expressed over the same scalar
quantity and the seam between them is well defined.

\subsection{The scalar estimand}
\label{sec:estimand}

The quantity to be protected is the receiver's error against a named laboratory
realization of Coordinated Universal Time,
\begin{equation}
  x \;=\; t_{\mathrm{rx}} - \utck{},
  \label{eq:estimand}
\end{equation}
a \emph{scalar}: unlike position ARAIM, whose estimand is a vector state observed
through a geometry matrix with dilution structure, every time source here reports
the same one-dimensional offset. This single fact---made precise in
Section~\ref{sec:scalar}---is what degenerates the geometry and moves the PL from
a geometry-driven to a bias-driven quantity. Fixing $\utck{}$ (rather than
``the server clock'') as the reference makes the traceability bias a first-class,
named term; its construction is deferred to a companion paper, and it enters here
only as an input overbound (Section~\ref{sec:sourcemodel}).

\subsection{The source model}
\label{sec:sourcemodel}

Each of $N$ heterogeneous time sources---representatively a GNSS precise-point
positioning source, a PTP/White-Rabbit packet source, and a low-Earth-orbit
signals-of-opportunity (LEO-STL) source---is reduced to its scalar $\utck{}$
contribution and summarized by three numbers:
\begin{equation}
  \text{source } i:\qquad
  \big(\,\sigma_i,\; b_i,\; p_{\mathrm{fault},i}\,\big),
  \label{eq:source}
\end{equation}
where $\sigma_i$ is a one-sigma Gaussian noise overbound on the source's $\utck{}$
estimate (seconds), $b_i$ is a \emph{bounded deterministic} $\utck{}$ traceability
bias overbound (seconds), and $p_{\mathrm{fault},i}$ is the source's prior fault
probability over the exposure interval. Two modelling commitments are load-bearing
and stated now.
\begin{itemize}
  \item \textbf{The bias $b_i$ is an input, not a claim of this paper.} It is a
  bounded deterministic overbound carried structurally separate from the Gaussian
  $\sigma_i$ tail: we never inflate a Type-B $k=2$ metrological uncertainty into a
  deep-tail Gaussian sigma (a category error), but treat the systematic component
  as a one-sided bounded bias in the ARAIM manner. Assembling a per-source library
  of $b_i$ from Circular-T $[\utck{}-\utck{}]$ and CGGTTS delay decompositions
  \citep{bipmcirculart,defraigne2015,panfilo2019}, with correlated
  cross-covariance for sources sharing a $\utck{}$, is the subject of a companion
  paper and is out of scope here; in this paper $b_i$ is simply given.
  \item \textbf{Independence is assumed within this paper and flagged as a
  limit.} The rank-one budget of Section~\ref{sec:scalar} treats the source fault
  modes as independent. A common-mode fault shared across sources (a shared
  $\utck{}$ realization; a GNSS system feeding two nominally independent sources)
  violates this and is not corrected here; the generalized-least-squares treatment
  of that case is a companion paper. We state the assumption rather than hide it.
\end{itemize}

\subsection{The handover state and the $\Pzero$ seed}
\label{sec:handover}

The disciplined oscillator that continues to keep time after the sources are lost
is modelled as a three-state clock---phase, fractional frequency, and
frequency drift---driven by the standard white-FM, random-walk-FM, and
random-run PSDs $(q_{\mathrm{wf}},q_{\mathrm{rw}},q_{\mathrm{drift}})$
\citep{zucca2005,riley2008,ieee1139}. At the instant of handover the disciplining
filter has a state-error covariance; its \emph{phase} block, the variance
$\Pzero$ (seconds$^2$), is the initial condition of the coast. This seed is not a
convenience: Section~\ref{sec:composition} proves that setting $\Pzero=0$ (coasting
as though the clock were perfectly known at handover) makes the holdover PL start
at zero and thereby \emph{drop} below the TPL that was being reported---the unsafe
seam. The coast additionally carries a deterministic aging term
$\tfrac12 D\tau^{2}$ (with $D$ the drift rate, $1/\mathrm{s}^2$) and a bounded
flicker-floor inflation; these are made precise in Section~\ref{sec:holdover}.
Throughout, phase-error variance grows as $\tau,\tau^{3},\tau^{5}$ for the
white-FM, random-walk-FM, and drift components respectively---the \emph{phase}
exponents, not the frequency-error exponents. The $\tau$ and $\tau^{3}$ terms are
the white-FM and random-walk-FM phase variances of the Zucca--Tavella model
\citep{zucca2005}; the $\tau^{5}$ term is the standard random-run (drift) extension
of that model \citep{riley2008,ieee1139}. We carry the phase-versus-frequency
distinction explicitly because conflating them is a common and consequential error.
